\documentclass[10pt]{article}
\usepackage[utf8]{inputenc}
\usepackage[T1]{fontenc}
\usepackage{amsmath}
\usepackage{amsfonts}
\usepackage{amssymb}
\usepackage[version=4]{mhchem}
\usepackage{stmaryrd}
\usepackage{bbold}
\usepackage{hyperref}
\hypersetup{colorlinks=true, linkcolor=blue, filecolor=magenta, urlcolor=cyan,}
\urlstyle{same}

\begin{document}
\section*{The Mathematics of Radical Denesting: Classical Theory, Galois Structures, and Computational Algorithms}
The Epistemological and Computational Significance of Nested Radicals

In the disciplines of algebraic number theory and symbolic computation, the representation and simplification of mathematical expressions are foundational to algorithmic efficiency and theoretical rigor. A nested radical is defined mathematically as an algebraic expression wherein a radical sign (such as a square root, cube root, or higher-order root) encapsulates another radical expression within its radicand ${ }^{1}$. Classic examples include expressions derived from regular polygons, geometric constructability problems, and the solutions to cubic equations, such as $\sqrt{5+2 \sqrt{6}}$ or the more complex configurations posited by Srinivasa Ramanujan, like $\sqrt[3]{\sqrt[3]{2}-1}{ }_{2}^{2}$. The algebraic operation of transforming these nested structures into equivalent forms with a reduced "nesting depth"-ideally removing the internal nesting entirely-is known as denesting ${ }^{2}$.\\
The mathematical imperative to denest radicals is not merely an aesthetic pursuit; it is a critical requirement for Computer Algebra Systems (CAS) and theoretical algorithms that process exact geometric computations. In computational mathematics, a persistent challenge is the "zero-equivalence problem," which requires algorithms to determine if a given complex algebraic expression evaluates identically to zero ${ }^{5}$. When dealing with radical expressions, denesting is often a prerequisite for establishing a canonical form ${ }^{6}$. Without a canonical representation, determining the linear independence of terms or establishing equality between two distinct formulations of the same algebraic number becomes computationally intractable, frequently leading to failures in differential calculus operations, numerical verifications, and geometric constraint logic programming ${ }^{3}$.\\
Historically, the simplification of nested radicals was relegated to heuristic pattern matching and ad-hoc algebraic manipulation. For decades, it was unclear whether a generalized algorithmic solution existed for radicals of arbitrary depth and degree. This uncertainty prevailed until the late 1980s and early 1990s, when Susan Landau, leveraging the principles of Galois theory and the Lenstra-Lenstra-Lovász (LLL) polynomial factorization algorithms, introduced the first deterministic algorithm capable of resolving the general denesting problem ${ }^{2}$. While Landau's work provided absolute theoretical closure by establishing necessary and sufficient conditions for denestability, the computational reality of its exponential time complexity spurred subsequent decades of research into bounded-degree algorithms, randomized Monte-Carlo zero-testing, and applied heuristic integrations within modern\\
software architectures ${ }^{5}$.

\section*{Classical Formulations and the Biquadratic Paradigm}
The foundational stratum of denesting theory focuses on the simplest nontrivial case: the depth-two nested square root involving strictly rational sub-components. The complete resolution of this biquadratic form represents a classical theorem that dictates precisely when and how a nested square root can be expanded into a linear combination of unnested square roots ${ }^{1}$.\\
For an algebraic expression of the form $\sqrt{a \pm \sqrt{c}}$, where both $a$ and $c$ are rational numbers and $c$ is not the square of a rational number, denesting is algebraically permissible if and only if $a$ is positive and the quantity $a^{2}-c$ is equal to the square of a positive rational number $d_{1}$. The quantity $a^{2}-c$ functions analogously to a discriminant in this context. When this stringent condition is met, the nested radical can be universally rewritten using the following identities:

$$
\begin{aligned}
& \sqrt{a+\sqrt{c}}=\sqrt{\frac{a+d}{2}}+\sqrt{\frac{a-d}{2}} \\
& \sqrt{a-\sqrt{c}}=\sqrt{\frac{a+d}{2}}-\sqrt{\frac{a-d}{2}}
\end{aligned}
$$

The algebraic proof of this theorem is remarkably elegant and relies on elementary field extensions and Vieta's formulas. Assuming that the nested radical can be expressed as a sum of two distinct square roots, one posits the equality $\sqrt{a \pm \sqrt{c}}=\sqrt{x} \pm \sqrt{y} \quad{ }^{1}$. By squaring both sides of this equation, the result is $a \pm \sqrt{c}=x+y \pm 2 \sqrt{x y}$. Because this equality must hold within the quadratic field extension $\mathbb{Q}(\sqrt{c})$, the rational and irrational components of both sides must be strictly equivalent. This forces the system of equations where $x+y=a \quad$ and $\quad 4 x y=c_{1}$.

According to Vieta's formulas, the variables $x$ and $y$ must therefore be the precise roots of the quadratic equation $X^{2}-a X+\frac{c}{4}=0$. Solving this quadratic equation yields the roots ${ }^{x}$ and ${ }^{y}$. For these roots to be valid rational numbers (which allows them to sit comfortably under the single square root in the denested form), the discriminant of the quadratic equation-which simplifies to $a^{2}-c$-must be a perfect square, and $a$ must be positive to ensure the roots themselves are positive real numbers ${ }^{1}$. If these conditions fail, the\\
expression cannot be denested into a simple sum or difference of square roots over the rational numbers.

\section*{Taxonomies of Denesting: Direct versus Indirect Modalities}
In contemporary pedagogical and theoretical frameworks, such as the 2017 formulations by Eleftherios Gkioulekas, the simplification of nested square roots is classified into two distinct operational modalities: direct denesting and indirect denesting ${ }^{11}$. This classification clarifies the boundaries of algebraic simplification when standard discriminant tests fail.\\
Direct denesting refers to the aforementioned classical scenario where the nested expression is rewritten exclusively as a sum or difference of square roots of rational numbers. The output remains structurally homogeneous, relying solely on index-two radicals ${ }^{11}$.\\
Indirect denesting addresses scenarios where an expression cannot be simplified into pure square roots, but can be factored by extracting a higher-order root of a rational number, leaving a simplified sum of square roots behind. Gkioulekas demonstrates that, based on earlier bounded-depth research by Borodin et al., if a nested square root expression requires higher-order roots for denesting, it will uniquely require fourth-order roots ${ }^{9}$. In an indirect denesting, the expression reduces to the form $\sqrt[4]{p}(\sqrt{x} \pm \sqrt{y})$.\\
The mathematical conditions governing these classifications highlight the strict boundaries of field extensions.

\begin{center}
\begin{tabular}[t]{|l|l|l|}
\hline
Denesting Classification & Target Algebraic Form & Requisite Condition (a,b,c,p $\in$ Q) \\
\hline
Direct Denesting & $\sqrt{x} \pm \sqrt{y}$ & $a^{2}-c=d^{2}$,where $d$ is a positive rational square. \\
\hline
Indirect Denesting & $\sqrt[4]{p}(\sqrt{x} \pm \sqrt{y})$ & $a^{2}-b^{2} p=\delta^{2} p$, where higher-degree roots facilitate the factoring of the minimal polynomial. \\
\hline
\end{tabular}
\end{center}

Galois theory further restricts the anatomy of these denestings. It is mathematically impossible to denest a two-term quadratic nested radical into a sum of three or more unnested radicals ${ }^{2}$.\\
Every element within the quadratic field $\mathbb{Q}(\sqrt{c})$ can be written uniquely as $\alpha+\beta \sqrt{c}$. Because field automorphisms must preserve the rational components while conjugating the irrational components, attempting to construct a denesting formula with three independent square roots invariably forces one or more coefficients to collapse to zero to maintain the symmetry of the Galois group. Consequently, the assumption that a generalized denesting requires a larger polynomial basis inevitably reduces back to the two-term format ${ }^{1}$.

\section*{Galois Theory and the Structural Homomorphisms of Subfields}
To progress beyond biquadratic expressions and address the denesting of radicals of arbitrary depth, indices, and degrees, the analytical framework must pivot from elementary algebra to Galois theory. Galois theory models the symmetries of the roots of polynomials, providing a definitive map of the subfield structures within an algebraic extension. The capability to denest a radical expression is fundamentally dictated by whether the algebraic number in question resides within a simpler subfield of its splitting field ${ }^{3}$.

When a mathematician or a computer algorithm encounters an algebraic number ${ }^{\alpha}$ defined by a nested radical, that number represents a root of a minimal polynomial $m_{\alpha}(x)$ over a designated base field $k$, which is traditionally the field of rational numbers $\mathbb{Q}_{8}$. The process of adjoining this root to the base field creates the algebraic extension $k(\alpha)$. In this context, denesting is not merely a simplification of notation; it is the rigorous search for an alternative, lower-degree generator for the same field, or the identification that ${ }^{\alpha}$ lies within an intermediate field extension that requires fewer radical operations to construct ${ }^{14}$. The mechanics of this are best illustrated by returning to the classic example $\alpha=\sqrt{5+2 \sqrt{6}}$. The algebraic number $\alpha$ satisfies the irreducible minimal polynomial $x^{4}-10 x^{2}+1=0$ over $\mathbb{Q}_{14}$. Adjoining ${ }^{\alpha}$ generates a field extension $\mathbb{Q}(\alpha)$ of degree 4 over the rationals. Using a standard power basis, any element in this field can be expressed as a linear combination of the set $\left\{1, \alpha, \alpha^{2}, \alpha^{3}\right\}^{14}$. However, because $\alpha$ simplifies to $\sqrt{2}+\sqrt{3}$, and because the square roots of 2, 3, and 6 are linearly independent over ${ }^{\mathbb{Q}}$, an alternative basis for the exact same field extension is $\{1, \sqrt{2}, \sqrt{3}, \sqrt{6}\}{ }^{14}$. Galois theory explains this equivalence through subfield structures. The field $\mathbb{Q}(\sqrt{2}, \sqrt{3})$ has a specific Galois group-the Klein four-group-which contains three nontrivial subgroups of index 2. By the Fundamental Theorem of Galois Theory, these subgroups correspond to exactly three intermediate quadratic subfields: $\mathbb{Q}(\sqrt{2}), \mathbb{Q}(\sqrt{3})$, and $\mathbb{Q}(\sqrt{6})^{14}$. The nested radical exists seamlessly within this architecture. Therefore, denesting algorithms essentially perform matrix transformations, mapping the complex nested power basis into the structurally simpler unnested basis defined by the subfields ${ }^{14}$.

\section*{Splitting Fields and Roots of Unity}
The theoretical boundaries of generalized denesting are governed by the properties of splitting fields and the presence of roots of unity. The roots of the minimal polynomial $m_{\alpha}(x)$ collectively generate a splitting field $L$ over the base field $k$. According to theorems derived from Artin and expanded by modern computer scientists, if the base field ${ }^{k}$ (assuming characteristic 0 ) already contains all requisite roots of unity, then a minimal-depth nesting of $\alpha$ exists such that every intermediate term lies strictly within the splitting field $L_{8}$. Complications arise because the rational field ${ }^{\mathbb{Q}}$ does not natively contain complex roots of unity (other than ${ }^{ \pm 1}$ ). When the base field lacks these roots, one must analyze the Galois group $G$ of the splitting field $L$ over $k$. By examining the derived series of $G$ and identifying the least common multiple $l$ of the exponents of this series, Galois theory proves that optimal denesting might be impossible strictly over $\mathbb{Q}$. Instead, the field must be artificially expanded by adjoining the ${ }^{l}$-th primitive root of unity, denoted $\zeta_{l 8}$. Adjoining this root of unity acts as an algebraic key, unlocking the cyclic extensions dictated by the Galois group and allowing the nested roots to untangle. If a theoretical denesting of depth $t_{\text {exists for an expression, there is a guaranteed denesting over the expanded field }} k\left(\zeta_{l}\right)$ with a depth of at most $t+1_{3}$. To restore absolute optimality and find the minimum possible depth, one must calculate the least common multiple $m$ of the indices of the roots in the original expression, and then adjoin a new root of unity $\zeta_{r}$, where $r^{r}$ is the least common multiple of $m$ and $l_{14}$. This Galois-centric approach forms the theoretical engine for all deterministic denesting algorithms.

\section*{The Advent of Deterministic Denesting: Landau's Algorithm}
Prior to 1989, the landscape of radical simplification was highly fragmented. The algorithms utilized by early computer algebra systems, such as MACSYMA, relied upon the foundational work of Caviness and Fateman (1976), which applied heuristic algebraic simplifications but lacked a unified theory capable of resolving arbitrary nestings ${ }^{3}$. Richard Zippel introduced a probabilistic algorithm in 1985 aimed at simplification, but his proofs regarding sufficient conditions for denestability contained subtle mathematical lacunae ${ }^{8}$. The general problem-whether an algorithm could deterministically decide if any given nested radical could\\
be denested-remained an open question in computer science ${ }^{2}$.\\
This computational void was definitively filled by Susan Landau. In a sequence of landmark papers culminating in her 1992 publication in the SIAM Journal on Computing, Landau introduced the first comprehensive algorithm for deciding denestability and actively computing the denested forms ${ }^{2}$. Landau's work corrected the gaps in Zippel's proofs, demonstrating that his sufficient conditions were also strictly necessary, thereby establishing a complete theoretical framework for the problem'.\\
Landau's algorithm operates deterministically by translating the algebraic expression into field theory mechanics. It takes a nested radical ${ }^{\alpha}$ and computes the splitting field of its minimal polynomial over ${ }^{k}$. The algorithm checks the base field for roots of unity; if they are present, it computes the exact optimal denesting. If they are absent, the algorithm calculates the necessary primitive root of unity $\zeta_{l}$, adjoins it to the field, and computes a denesting that is within depth one of the absolute mathematical optimum, utilizing the symbol $\zeta_{l}$ to represent the root rather than further nesting the expression ${ }^{3}$.

\section*{The Exponential Time Complexity Bottleneck}
While Landau's algorithm represented a monumental triumph in abstract mathematics, its translation into applied computer science revealed a severe computational limitation. The algorithm relies inextricably on computing the splitting field of the minimal polynomial of $\alpha$.

The degree of a splitting field is highly volatile; for a polynomial of degree $n$, the splitting field can, in the worst-case scenario, reach a degree of $n!$ (n-factorial) ${ }^{14}$. Consequently, Landau's algorithm possesses an exponential time complexity relative to the depth and degree of the nested radicals ${ }^{1}$.\\
To execute the necessary polynomial factorizations over these massive algebraic number fields, the algorithm utilizes advanced lattice basis reduction techniques, specifically the Lenstra-Lenstra-Lovász (LLL) algorithm ${ }^{8}$. Although the LLL algorithm functions in polynomial time for factoring over rational numbers, the sheer size of the splitting fields required for deeply nested radicals renders the overall denesting operation intractable for real-time computing environments ${ }^{1}$. As later literature and algorithmic research indicate, finding a deterministic algorithm that operates in polynomial time for the general denesting problem remains a significant open challenge ${ }^{10}$. Because of this exponential bottleneck, Landau's algorithm serves more as an existence proof and a theoretical benchmark rather than a routine subroutine in consumer software.

\section*{Bounded Degrees, Cubic Radicals, and the Ramanujan Identities}
Given the exponential constraints of generalized denesting via splitting fields, computational\\
mathematicians pivoted toward optimizing specific, bounded subsets of the problem. This shift in focus was profoundly influenced by the historical notebooks of Srinivasa Ramanujan. Ramanujan possessed an extraordinary intuition for discovering highly complex, non-obvious algebraic identities involving cubic and higher-order nested radicals ${ }^{16}$. Some of his most famous finite denesting identities include:

$$
\begin{gathered}
\sqrt[3]{\sqrt[3]{2}-1}=\sqrt[3]{\frac{1}{9}}-\sqrt[3]{\frac{2}{9}}+\sqrt[3]{\frac{4}{9}} \\
\sqrt{\sqrt[3]{5}-\sqrt[3]{4}}=\frac{1}{3}(\sqrt[3]{2}+\sqrt[3]{20}-\sqrt[3]{25}) \\
\sqrt[6]{7 \sqrt[3]{20}-19}=\sqrt[3]{\frac{5}{3}}-\sqrt[3]{\frac{2}{3}}
\end{gathered}
$$

These identities are mathematically jarring because the left-hand expressions feature depth-two nestings of cubic and sixth roots, while the right-hand expressions collapse into depth-one linear combinations ${ }^{2}$. For decades, these stood as isolated mathematical curiosities. However, in 1992, Johannes Blömer published "How to Denest Ramanujan's Nested Radicals," which provided the formal algorithmic structure explaining why these specific identities work, without resorting to the exponential splitting fields required by Landau's method ${ }^{9}$. Blömer's breakthrough focused on optimizing depth-two nested radicals involving bounded degrees. To denest an expression, an algorithm must first determine if a given sum of radicals evaluates to zero. Blömer achieved this by transforming the sum into a basis of linearly independent radicals using a theorem originally posited by C.L. Siegel ${ }^{5}$. Siegel's Theorem dictates that for any real field, if the ratio of any two distinct radicals in a set does not reside within the base field, then the entire set of radicals is linearly independent ${ }^{5}$. By substituting exponential field factorizations with a streamlined ratio test, Blömer developed a randomized Monte-Carlo algorithm that operates in polynomial time relative to $\log \left(d_{i}\right)$ (where $d_{i}$ are the degrees of the roots) ${ }^{5}$. This provided a near-optimal depth denesting for bounded degree radicals and fully demystified the algebraic structure underlying Ramanujan's cubic anomalies.

\section*{Modern Factorization Equivalences in Cubic Denesting}
Research into cubic denesting continues to yield sophisticated theoretical insights. A 2024 mathematical paper by Alberto Cavallo, "Denesting cubic radicals," provides an exhaustive analysis of the conditions required to simplify radicals of the form $\sqrt[3]{a+b \sqrt{p}}$, providing both a complete theoretical answer and a corresponding programmatic implementation ${ }^{4}$. Cavallo's framework assumes the cubic radical can be denested into a simpler form $A+B \sqrt{p}$, where $A, B \in \mathbb{Q}$. This implies that the field $\mathbb{Q}(\sqrt[3]{a+b \sqrt{p}})$ functions as a 2-dimensional vector space over the rational numbers ${ }^{4}$. By cubing the assumed identity, one\\
derives a rigid system of equations: $a=A\left(A^{2}+3 B^{2} p\right) \quad$ and $b=B\left(3 A^{2}+B^{2} p\right)$. The norm of this extension is defined as $N=a^{2}-b^{2} p$, allowing the minimal polynomial of the radical to be deduced ${ }^{4}$.\\
The critical innovation in Cavallo's work is the reduction of the denesting problem to the identification of rational roots of a depressed cubic polynomial $h_{1}(x)=x^{3}+P x+Q$. By mapping the denesting constraints onto this cubic polynomial, Cavallo proves that determining whether a rational cubic polynomial is reducible-and thus, whether the nested radical can be simplified-is algorithmically equivalent in complexity to finding the prime factorization of natural numbers ${ }^{4}$. This establishes a rigorous computational boundary: cubic denesting is exactly as computationally hard as integer factorization, linking the algebraic simplification of radicals to foundational concepts in cryptography and prime number theory.

\section*{The Analytic Divergence: Infinitely Nested Radicals}
While the denesting of finite radicals is strictly a problem of abstract algebra and field extensions, extending the nesting process to infinity shifts the mathematical paradigm entirely into the domain of real analysis, calculus, and limit convergence ${ }^{16}$. Infinitely nested radicals (also known as continued roots) require continuity proofs and functional analysis to evaluate.\\
This intersection of domains was famously highlighted by Srinivasa Ramanujan in 1911 when he submitted a seemingly impossible problem to the Journal of the Indian Mathematical Society ${ }^{16}$. He asked readers to evaluate the infinite sequence:

$$
\sqrt{1+2 \sqrt{1+3 \sqrt{1+4 \sqrt{1+\cdots}}}}
$$

When no solutions were provided by the mathematical community, Ramanujan published his own resolution, bypassing field extensions entirely in favor of functional algebra ${ }^{16}$. He defined a generalized identity based on polynomial expansion:

$$
x+n+a=\sqrt{a x+(n+a)^{2}+x \sqrt{a(x+n)+(n+a)^{2}+(x+n) \sqrt{\cdots}}}
$$

By meticulously substituting $x=2, n=1$, and $a=0$ into this infinite sequence, the right-hand side exactly matches the journal prompt, and the left-hand side evaluates simply to $2+1+0=3{ }^{16}$. Modern calculus verifies this through recursive functions. By setting a function $f(x)=x+1$, one observes that squaring the function yields $(x+1)^{2}=x^{2}+2 x+1$, which can be factored into $1+x(x+2)$, equivalent to $1+x f(x+1)$. Taking the square root of both sides creates the recursive definition\\
$f(x)=\sqrt{1+x f(x+1)}$, which endlessly expands into Ramanujan's nested sequence ${ }^{16}$.\\
The convergence of such infinite radicals is governed by strict analytic rules, most notably Herschfeld's Convergence Theorem. Herschfeld proved that an infinitely nested radical of the form $\sqrt{a_{1}+\sqrt{a_{2}+\ldots}}$ (where all $a_{i}$ are positive real numbers) converges to a finite limit if and only if there exists a bounding constant $M$ such that $a_{n}<M^{2^{n}}$ for all integers $n_{1}$. When the terms are strictly repeating, the limits frequently converge on foundational mathematical constants. For example, the infinitely repeating radical $y=\sqrt{1+\sqrt{1+\sqrt{1+\cdots}}} \quad$ is evaluated by recognizing its self-similarity as $y=\sqrt{1+y}$. Squaring this yields the characteristic quadratic equation $y^{2}-y-1=0$, the positive root of which is the Golden Ratio, $\phi=\frac{1+\sqrt{5}}{2}{ }^{1}$. Similarly, the infinite nesting of subtraction, $y=\sqrt{n-\sqrt{n-\cdots}}$, resolves to the positive root of $y^{2}+y-n=0_{1}^{1}$. The concept also intersects with trigonometry and geometry, most famously in Viète's historical formula for ${ }^{\pi}$. Discovered in the 16th century, Viète's formula utilizes an infinite product of nested square roots of 2 to approximate the area of a circle, standing as one of the earliest examples of an infinite product in mathematical analysis ${ }^{1}$ :

$$
\frac{2}{\pi}=\frac{\sqrt{2}}{2} \cdot \frac{\sqrt{2+\sqrt{2}}}{2} \cdot \frac{\sqrt{2+\sqrt{2+\sqrt{2}}}}{2} \cdots
$$

\section*{Applied Computational Algorithms in Modern Software}
The translation of radical denesting theories from the chalkboard to commercial and open-source Computer Algebra Systems highlights a persistent tension between mathematical completeness and computational pragmatism. While Landau proved that general denesting is deterministically solvable, the exponential time complexity associated with computing minimal polynomials and splitting fields over Galois groups renders her exact algorithm unfeasible for consumer-facing software ${ }^{2}$. Users of CAS expect near-instantaneous simplifications; an algorithm that locks the CPU into calculating an $n!$-degree lattice reduction is practically useless for routine engineering or physical modeling tasks. Consequently, system developers rely on a mosaic of heuristic algorithms, recursive pattern matching, and bounded-depth logic ${ }^{23}$.

\begin{center}
\begin{tabular}[t]{|l|l|l|}
\hline
Computer Algebra System & Primary Denesting Module & Algorithmic Implementation Strategy \\
\hline
Maxima & sqdnst.mac / raddenest & Originally relied on brutal unguided decomposition. Patched by community members to implement the Jeffrey/Rich recursive algorithms and Borodin's depth-reduction techniques ${ }^{23}$. \\
\hline
SymPy (Python) & sympy/simplify/sqrtdenest.p y & Utilizes a recursive, modular testing approach. Checks sub-discriminants iteratively. Built for speed and stability rather than theoretical Galois completeness ${ }^{23}$. \\
\hline
Maple & radnormal & Operates aggressively in the background, utilizing radical normal forms to force standard representations of inputs to prevent zero-equivalence errors in later calculus operations ${ }^{6}$. \\
\hline
\end{tabular}
\end{center}

\section*{Implementations in Maxima and SymPy}
In the open-source CAS Maxima, the core function sqdnst.mac was designed around the Jeffrey and Rich heuristic identity, which algebraically forces the expression into the form: $\sqrt{X+Y}=\sqrt{X / 2+\sqrt{X^{2}-Y^{2}} / 2}+\sqrt{X / 2-\sqrt{X^{2}-Y^{2}} / 2} \quad$ 23. While theoretically sound, early implementations executed this strategy poorly. If presented with a multi-term sum under a radical, the system would arbitrarily split the first term from the rest and attempt the formula. If the resulting discriminant $\sqrt{X^{2}-Y^{2}}$ was not a perfect square, the denesting failed. The system did not recursively check alternate decompositions ${ }^{23}$. To remedy this, the Maxima community introduced raddenest, a module explicitly built to incorporate Borodin's depth-decreasing algorithms. While vastly more capable at handling deeply nested terms, developers acknowledge that it still falls short of the absolute generality\\
provided by Zippel or Landau ${ }^{23}$.\\
SymPy, the primary symbolic mathematics library for Python, handles radical simplification through the \href{http://sqrtdenest.py}{sqrtdenest.py} module ${ }^{23}$. Because Python is an interpreted language and inherently lacks the low-level execution speed required for massive algebraic lattice reductions, SymPy developers have consciously avoided implementing Landau's splitting field model ${ }^{23}$. Instead, \href{http://sqrtdenest.py}{sqrtdenest.py} focuses on decomposing radical additions and rigorously checking for perfect squares within sub-expressions ${ }^{25}$. To maintain stability, the module is backed by an extensive suite of unit tests, designed to catch specific edge cases and prevent the recursive loops that often plague heuristic simplifiers ${ }^{28}$. When SymPy encounters a highly contrived, deeply nested radical that defies its heuristic checks, the software defaults to safety: it returns the unsimplified nested radical rather than risk a timeout error or an incorrect evaluation.

\section*{Maple and the Radical Normal Form}
Proprietary systems like Maple often take a more aggressive stance on simplification. Maple utilizes the radnormal command, which seeks to immediately establish a "radical normal form" for inputted expressions ${ }^{6}$. Maple's philosophy is preemptive; by aggressively attempting to denest and normalize radicals as soon as they are entered, the software prevents the accumulation of unsimplified terms that cause zero-testing failures during definite integration or Hermite reduction processes ${ }^{6}$. While Maple's exact proprietary codebase is closed, its ability to parse hypergeometric functions and complex bounded-degree radicals suggests an integration of Blömer's bounded zero-equivalence models alongside traditional heuristic mapping ${ }^{5}$.

\section*{Conclusion}
The study of radical denesting sits at a remarkable crossroads between elementary arithmetic, abstract algebra, and computational complexity. The progression of the field illustrates a continuous effort to bring the elegant symmetries of Galois theory into the constrained reality of computer architecture. The classical resolution of the biquadratic form $\sqrt{a \pm \sqrt{c}}$ provides a pristine example of discriminant logic and field extensions operating in perfect harmony. Yet, as the depth and degree of the radicals increase, this algebraic simplicity rapidly escalates into topological complexity, requiring the adjunction of roots of unity and the exhaustive mapping of polynomial splitting fields.\\
Susan Landau's deterministic algorithms represent the theoretical pinnacle of this discipline, providing the necessary mathematical closure that proves general denesting is intrinsically solvable. However, the exponential time complexity inherent in calculating Galois structures dictates that Landau's algorithm remains a mathematical ideal rather than a computational standard. Modern researchers and software engineers have adapted to this reality by isolating bounded degrees and specific cubic structures-such as those famously intuited by Srinivasa Ramanujan-proving that subsets of the denesting problem can be solved in polynomial time through Monte-Carlo testing and connections to prime factorization.\\
Ultimately, modern Computer Algebra Systems reflect a pragmatic compromise. By leveraging heuristic identities, recursive pattern matching, and radical normal forms, software like SymPy,

Maxima, and Maple efficiently navigate the vast majority of real-world geometric and physical computing tasks. As hardware capabilities expand and algorithms for lattice basis reductions optimize, the gap between Landau's abstract Galois mapping and applied computational simplification will continue to narrow, further untangling the nested structures of algebraic numbers.

\section*{Works cited}
\begin{itemize}
  \item[1.] Nested radical - Wikipedia, \href{https://en.wikipedia.org/wiki/Nested_radical}{https://en.wikipedia.org/wiki/Nested\_radical}
  \item[2.] Understanding Nested Radicals in Algebra | PDF - Scribd, \href{https://www.scribd.com/document/968077385/Nested-Radical-Wikipedia-1}{https://www.scribd.com/document/968077385/Nested-Radical-Wikipedia-1}
  \item[3.] (PDF) Simplification of Nested Radicals - ResearchGate, \href{https://www.researchgate.net/publication/2629046_Simplification_of_Nested_Rad}{https://www.researchgate.net/publication/2629046\_Simplification\_of\_Nested\_Rad} icals
  \item[4.] Denesting cubic radicals - arXiv, \href{https://arxiv.org/pdf/2403.04776}{https://arxiv.org/pdf/2403.04776}
  \item[5.] Simplifying Expressions Involving Radicals - NYU Computer Science, \href{https://cs.nyu.edu/exact/pap/rootBounds/sumOfSqrts/bloemerThesis.pdf}{https://cs.nyu.edu/exact/pap/rootBounds/sumOfSqrts/bloemerThesis.pdf}
  \item[6.] The Didactical Challenge of Symbolic Calculators - dokumen.pub, \href{https://dokumen.pub/download/didactical-challenge-of-symbolic-calculators-tur}{https://dokumen.pub/download/didactical-challenge-of-symbolic-calculators-tur} ning-a-computanal-device-into-a-mathematical-instrument-the.html
  \item[7.] nth root - Wikipedia, \href{https://en.wikipedia.org/wiki/Nth_root}{https://en.wikipedia.org/wiki/Nth\_root}
  \item[8.] Simplification of Nested Radicals | SIAM Journal on Computing, \href{https://epubs.siam.org/doi/10.1137/0221009}{https://epubs.siam.org/doi/10.1137/0221009}
  \item[9.] [PDF] Simplification of nested radicals - Semantic Scholar, \href{https://www.semanticscholar.org/paper/Simplification-of-nested-radicals-Landau}{https://www.semanticscholar.org/paper/Simplification-of-nested-radicals-Landau} /37b2e1949be7d1f950e05bcfdcfb31daa00a970a
  \item[10.] Relevance of Landau's Algorithm for Denesting Radicals, \href{https://mathoverflow.net/questions/319874/relevance-of-landaus-algorithm-for-d}{https://mathoverflow.net/questions/319874/relevance-of-landaus-algorithm-for-d} enesting-radicals
  \item[11.] On the denesting of nested square roots - UTRGV Faculty Web, \href{https://faculty.utrgv.edu/eleftherios.gkioulekas/papers/submitted/denesting-squar}{https://faculty.utrgv.edu/eleftherios.gkioulekas/papers/submitted/denesting-squar} e-roots.pdf
  \item[12.] International Journal of Mathematical Education in Science and, \href{https://www.tandfonline.com/toc/tmes20/48/6}{https://www.tandfonline.com/toc/tmes20/48/6}
  \item[13.] A Classical Introduction to Galois Theory | Wiley, \href{https://www.wiley.com/en-us/a-classical-introduction-to-galois-theory-p-9781118}{https://www.wiley.com/en-us/a-classical-introduction-to-galois-theory-p-9781118} 336816
  \item[14.] $2+\sqrt{ } 3$ : four different views - Susan Landau, \href{https://privacyink.org/pdf/Sqrt2+Sqrt3.pdf}{https://privacyink.org/pdf/Sqrt2+Sqrt3.pdf}
  \item[15.] A135611 - OEIS, \href{https://oeis.org/A135611/internal}{https://oeis.org/A135611/internal}
  \item[16.] Ramanujan's Nested Radical Problem | by Ujjwal Singh, \href{https://www.cantorsparadise.com/ramanujans-nested-radical-problem-f14c3060}{https://www.cantorsparadise.com/ramanujans-nested-radical-problem-f14c3060} e495
  \item[17.] Denesting Radicals in Ramanujan's Work | PDF | Field (Mathematics, \href{https://www.scribd.com/document/264470815/Ramanujan-Day}{https://www.scribd.com/document/264470815/Ramanujan-Day}
\end{itemize}

\begin{itemize}
  \item[18.] (PDF) Variations on Ramanujan's nested radicals - ResearchGate, \href{https://www.researchgate.net/publication/284579347}{https://www.researchgate.net/publication/284579347} Variations on Ramanujan's nested radicals
  \item[19.] [2403.04776] Denesting cubic radicals - arXiv, \href{https://arxiv.org/abs/2403.04776}{https://arxiv.org/abs/2403.04776}
  \item[20.] Denesting cubic radicals - arXiv, \href{https://arxiv.org/html/2403.04776v2}{https://arxiv.org/html/2403.04776v2}
  \item[21.] The Infinitely Nested Radicals Problem and Ramanujan's wondrous, \href{https://math.stackexchange.com/questions/4049946/the-infinitely-nested-radical}{https://math.stackexchange.com/questions/4049946/the-infinitely-nested-radical} s-problem-and-ramanujans-wondrous-formula
  \item[22.] Calculus of Ramanujan's Infinite Nested Radicals | by Archie Smith, \href{https://medium.com/@ArchieGSmith/calculus-of-infinite-nested-radicals-1ab3e1}{https://medium.com/@ArchieGSmith/calculus-of-infinite-nested-radicals-1ab3e1} d26774
  \item[23.] [Maxima-discuss] Denesting, \href{https://maxima-discuss.narkive.com/rz7bUs4t/denesting}{https://maxima-discuss.narkive.com/rz7bUs4t/denesting}
  \item[24.] [Maxima-discuss] Denesting radical expressions (package raddenest), \href{https://sourceforge.net/p/maxima/mailman/maxima-discuss/thread/oblan6$h4r$1}{https://sourceforge.net/p/maxima/mailman/maxima-discuss/thread/oblan6\$h4r\$1} @blaine.gmane.org/
  \item[25.] miniconda3/envs/poem/lib/python3.10/site-packages/sympy/simplify, \href{https://mo.zju.edu.cn/repo/OZ2fxZukDanlbwK8WP7Ht}{https://mo.zju.edu.cn/repo/OZ2fxZukDanlbwK8WP7Ht} iyYx-iFiohusxaM3QPIs2m/ blob/a8e02444aaf0d71155a42be0422841306d4d7d8a/miniconda3/envs/poem/lib /python3.10/site-packages/sympy/simplify/sqrtdenest.py
  \item[26.] Overview - Maple Help - Maplesoft, \href{https://jp.maplesoft.com/support/help/maple/view.aspx?path=updates/v53&L=J}{https://jp.maplesoft.com/support/help/maple/view.aspx?path=updates/v53\&L=J}
  \item[27.] port math/py-sympy - OpenBSD Ports Readme, \href{https://openports.pl/path/math/py-sympy}{https://openports.pl/path/math/py-sympy}
  \item[28.] Writing Tests - SymPy 1.14.0 documentation, \href{https://docs.sympy.org/latest/contributing/new-contributors-guide/writing-tests}{https://docs.sympy.org/latest/contributing/new-contributors-guide/writing-tests}. html
\end{itemize}

\end{document}